\subsubsection{Orthogonal quartic blocks and their weights}
\label{sec:Clifford-composition-background}

The OT--FKM and Riedler quartics are quadratic profiles in Clifford
coordinates and the radius. The spin-factor realization of
\cref{sec:quadratic-jordan-framework} gives their common quartic table.
The underlying Hurwitz matrix construction and its reduction to
Clifford representations are described in
\cite[Chapter~0, pp.~2--3; Chapter~1, Lemma~1.5 and Proposition~1.7]{Shapiro2000}.

Let $\mathcal C=(P_0,\ldots,P_m)$ be a symmetric Clifford system on
$\R^{2\ell}$, with $m\geq1$, and put
\[
 R=|z|^2,\qquad p_i(z)=\langle P_i z,z\rangle,
 \qquad p=(p_0,\ldots,p_m).
\]
Their differential table is \eqref{eq:Clifford-pullback-identities},
and \cref{lem:Clifford-algebra-pullback} gives $|p|^2\leq R^2$.

In the product $f\circ g=\Gamma(f,g)/4$, the quadratic space
$J=\operatorname{span}\{R,p_0,\ldots,p_m\}$ has multiplication
\begin{equation}\label{eq:Clifford-composition-quadratic-table}
 R\circ R=R,\qquad R\circ p_i=p_i,\qquad
 p_i\circ p_j=\delta_{ij}R.
\end{equation}
Thus $J$ is a spin-factor Jordan algebra. The Clifford relations
give the following table for its quartic block norms.

\begin{lemma}[Orthogonal quartic blocks]
\label{lem:Clifford-block-quartics}
Partition $\{0,\ldots,m\}$ into nonempty sets $I_1,\ldots,I_s$,
and set
\[
 E_j=\sum_{i\in I_j}p_i^2\quad(1\leq j\leq s),\qquad
 E_*=R^2-\sum_{i=0}^m p_i^2.
\]
Then all these polynomials are nonnegative and
\begin{align}
 &E_*+\sum_{j=1}^s E_j=R^2,\qquad
 \Gamma(E_i,E_j)=16RE_i\delta_{ij},
 \label{eq:Clifford-composition-quartic-table}\\
 &\Delta E_j=8|I_j|R,\qquad
 \Delta E_*=8(\ell-m)R.\nonumber
\end{align}
Here the pairing identity includes the index $*$.
\end{lemma}

\begin{proof}
For disjoint blocks the pairing in
\eqref{eq:Clifford-pullback-identities} gives
$\Gamma(E_i,E_j)=0$; for one block it gives
$\Gamma(E_j,E_j)=16RE_j$ and $\Delta E_j=8|I_j|R$.
Also $\Gamma(R^2,E_j)=16RE_j$, so $E_*$ is orthogonal to every
$E_j$.  Expanding its gradient square gives
$\Gamma(E_*,E_*)=16RE_*$.  Finally,
$\Delta R^2=8(\ell+1)R$ proves the asserted Laplacian of $E_*$.
\end{proof}

\begin{remark}[Geometric interpretation]
\label{rem:weighted-Clifford-foliation}
For $\ell>m+1$, the Clifford map on the unit sphere has connected
fibers and hemisphere quotient metric. If all $|I_j|\geq2$, pulling
back the target $\prod_j\mathrm{SO}(|I_j|)$-orbits gives the common
quartic-block levels: this is Radeschi's composed foliation
\cite[Theorems~A--B and Section~3]{Radeschi14}. The one-leaf target
gives the OT--FKM foliation. In the splitting model the target
$\mathrm{SO}(m)$-action fixes $p_0$; its leaves retain $u-v$, while
the quartic subalgebra retains only $(u-v)^2$.
The balanced $(3,7,7)$ model in
\cref{sec:weighted-Clifford-quartic} uses $\mathcal C_{9,1}$ and
target $\mathrm{SO}(3)\times\mathrm{SO}(7)$-orbits, with profile
$E_+=E_-$. A profile equation selects a union of leaves; the
differential table requires no connectedness assumption.
\end{remark}

For three quartic blocks, the weights in the table are the target
block dimensions and the remaining weight $\ell-m$. The following
criterion identifies every minimal linear combination with a regular
zero meeting the positive part of this table.

\begin{proposition}[Minimal levels of orthogonal quartic generators]
\label{prop:orthogonal-quartics}
Suppose that nonzero homogeneous quartic polynomials $E_0,E_1,E_2$ satisfy
\begin{equation}\label{eq:orthogonal-quartic-data}
 E_0+E_1+E_2=R^2,\qquad
 \Gamma(E_i,E_j)=16RE_i\delta_{ij},\qquad
 \Delta E_i=8r_iR
\end{equation}
for constants $r_i>0$. They generate a Laplacian algebra together
with $R$, of generator degrees $(2,4,4)$ after eliminating one $E_i$.
Let $F=\sum_i\alpha_iE_i$ have a regular zero point with all $E_i>0$,
and put $\sigma=\sum_i r_i\alpha_i$.
Then $F=0$ is minimal if and only if its nonzero coefficients have
exactly two values $\alpha>0>\beta$, satisfying
\begin{equation}\label{eq:quartic-level-condition}
 \alpha(r_+-1)+\beta(r_--1)=0,\qquad
 r_+=\sum_{\alpha_i=\alpha}r_i,\quad r_-=\sum_{\alpha_i=\beta}r_i.
\end{equation}
If $F=0$ is minimal and none of the coefficients vanishes,
$\R[R,F]$ is already Laplacian.
\end{proposition}

\begin{proof}
The norm identities give $E_i\geq0$. Where they are positive,
their gradient matrix is positive definite, so their image has
interior. Set $A_j=\sum_i\alpha_i^jE_i$. The closure table gives
\[
 \Gamma(F,F)=16RA_2,\qquad
 \cL(F)=128R^2(\sigma A_2-A_3)-64FA_2.
\]
Thus minimality on the regular zero plane is equivalent to
$\sigma A_2-A_3=\kappa F$ for a constant $\kappa$. The image contains
an open piece of that plane, so comparison of its linear coefficients
is valid. For each nonzero $\alpha_i$ this reads
$\alpha_i(\sigma-\alpha_i)=\kappa$. There are at most two such
values, and a zero point with positive $E_i$ requires both signs.
Their sum is $\sigma$, which is precisely
\eqref{eq:quartic-level-condition}. These conditions also make
$\cL(F)$ divisible by $F$ globally and prove sufficiency.
In the minimal case, if all coefficients are nonzero, then
$A_2=(\alpha+\beta)F-\alpha\beta R^2$; the displayed norm and
Laplacian identities prove closure of $\R[R,F]$.
\end{proof}

\begin{remark}[More quartic blocks]\label{rem:quartic-many-blocks}
For any number of nonzero orthogonal quartic blocks summing to $R^2$,
the same coefficient comparison permits only two nonzero coefficient
values, provided the regular zero meets the positive-block region.
Group equal positive, negative, and zero coefficients and add their
weights. This reduces the minimal-level equation to
\eqref{eq:quartic-level-condition}, with two generators when no block
is unused. Different representations can still yield different cones.
\end{remark}

\begin{remark}[Compatible Cartan--M\"unzner equations]
\label{rem:compatible-quartic-CM}
Setting $\Phi_i=R^2-2E_i$ in \eqref{eq:orthogonal-quartic-data} gives
\[
 \sum_i\Phi_i=R^2,\qquad |\nabla\Phi_i|^2=16R^3,\qquad
 \Delta\Phi_i=4(N+2-4r_i)R.
\]
For distinct $i,j,k$ it also gives
$\Gamma(\Phi_i,\Phi_j)=-16R\Phi_k$.
Conversely, the sum and the individual gradient identities imply
these cross pairings; together with the displayed Laplacians they
recover \eqref{eq:orthogonal-quartic-data}.
These are compatible Cartan--M\"unzner-type equations, allowing
degenerate quartic representations: for example,
$E_i=(|x|^2-|y|^2)^2$ with equal-dimensional blocks has $r_i=1$.
\end{remark}

For the quartic algebra $\cA=\R[R,E_0,E_1,E_2]$ itself,
$\cA_2=\R R$. In the Clifford construction above, it lies in a
larger quadratic algebra. The criterion below determines its
realizable weights.

\begin{lemma}[Weights and image of the quartic blocks]
\label{lem:quartic-block-weights-image}
For the nonzero blocks in \eqref{eq:orthogonal-quartic-data}, the
weights are positive integers and $N=2(r_0+r_1+r_2-1)$. Their
spherical image is exactly
\begin{equation}\label{eq:quartic-simplex-image}
 (E_0,E_1,E_2)(\Sph^{N-1})
 =\{(u_0,u_1,u_2)\in[0,1]^3:u_0+u_1+u_2=1\}.
\end{equation}
\end{lemma}

\begin{proof}
Summing the Laplacian identities in \eqref{eq:orthogonal-quartic-data}
gives $N=2(r_0+r_1+r_2-1)$. The norm identities imply $E_i\geq0$.
On the unit sphere,
$u_i=E_i|_{\Sph^{N-1}}$ satisfies
\[
 |\nabla_{\Sph}u_i|^2=16u_i(1-u_i),\qquad
 \Delta_{\Sph}u_i=8r_i-4(N+2)u_i.
\]
A constant $u_i$ would be zero or one, contradicting the assumption
that all three blocks are nonzero.  Its minimum and maximum are
therefore zero and one.  Integration of the Laplacian identity
gives $r_i>0$.  At a nonzero zero of $E_i$, write
$H_i=\Hess E_i$.  Comparing quadratic terms in
$|\nabla E_i|^2=16RE_i$ gives
\[
 H_i^2=8RH_i,\qquad \operatorname{tr}H_i=8r_iR.
\]
Thus $H_i=8RP_i$ for an orthogonal projection $P_i$ of rank $r_i$;
in particular, each weight is an integer.

Nonnegativity and the sum of the blocks place the spherical image
in the displayed simplex. Its intersection with the simplex interior
is nonempty, since $E_0E_1E_2$ is a nonzero polynomial. There the
spherical gradient matrix
$16(\operatorname{diag}(u)-uu^{\mathsf T})$ has rank two, so the
intersection is relatively open by the submersion theorem. It is
also relatively closed, since the spherical image is compact.
Connectedness of the simplex interior, followed by taking its closure,
proves \eqref{eq:quartic-simplex-image}.
\end{proof}

The Clifford-module restriction gives an exact numerical range for
\eqref{eq:orthogonal-quartic-data}. We use the conventional function
\[
 (\delta(1),\ldots,\delta(8))=(1,2,4,4,8,8,8,8),
 \qquad \delta(k+8)=16\delta(k).
\]
Here $\delta(k)$ is the real irreducible module dimension for the
negative Clifford algebra on $k-1$ generators. A symmetric Clifford
system with $p\ge2$ matrices on $\R^{2\ell}$ exists exactly when
$\delta(p-1)\mid\ell$ \cite[Section~3.5]{FKM81}.
See also \cite[Chapter~0, pp.~3--4]{Shapiro2000} for the equivalent
Hurwitz--Radon restriction.

\begin{proposition}[Realizable quartic weights]
\label{prop:quartic-weight-realizability}
Let $r_0,r_1,r_2$ be positive integers, write their rearrangement as
$a\le b\le c$, and set $\ell=a+b+c-1$.
Nonzero quartics satisfying \eqref{eq:orthogonal-quartic-data} on
$\R^{2\ell}$ exist if and only if \(\delta(a+b-1)\mid\ell\).
Whenever this condition holds, a realization is obtained from one
Clifford system by two blocks of squared quadratic coordinates and
their complement in $R^2$.
\end{proposition}

\begin{proof}[Construction and parameter restrictions]
\leavevmode\par\nobreak
Assume the divisibility condition. Take $a+b$ symmetric Clifford
matrices, with quadratic coordinates $p_1,\ldots,p_{a+b}$.
The three polynomials
\[
 \sum_{\alpha=1}^{a}p_\alpha^2,\qquad
 \sum_{\alpha=a+1}^{a+b}p_\alpha^2,\qquad
 R^2-\sum_{\alpha=1}^{a+b}p_\alpha^2
\]
have weights $(a,b,c)$ by \cref{lem:Clifford-block-quartics}.
Their positive Laplacian coefficients ensure that none is zero.

For any realization, summing the Laplacians forces
$N=2(r_0+r_1+r_2-1)$. Positivity and integrality of the weights follow from
\cref{lem:quartic-block-weights-image}.
The converse uses the completed classification with four principal
curvatures and compatibility of the three blocks. See
\cref{app:quartic-weight-realizability} and its weight-one,
Clifford-joining, and exceptional-triple lemmas
\cref{lem:quartic-weight-one,lem:quartic-Clifford-joining,lem:quartic-exclude-155}.
\end{proof}

This is a classification of realizable weights, not of simultaneous
orthogonal equivalence classes of triples. For example, $(3,3,5)$
is excluded since $\delta(5)=8\nmid10$, although the individual
Cartan--M\"unzner multiplicity pairs $(2,7),(2,7),(4,5)$ all occur.
The criterion is symmetric in the weights; the subcalibration
range in \cref{thm:quartic-table-existence} instead depends on
which two blocks are chosen for the defining polynomial.

\begin{corollary}[The weights with active sum nine]
\label{cor:quartic-nine-weights}
Let $r_0,r_1,r_2$ be positive integers with $r_1,r_2\geq2$ and
$r_1+r_2=9$. A quartic realization exists if and only if $r_0$ is
a positive multiple of eight. Its ambient dimension is $N=2r_0+16$.
\end{corollary}

\begin{proof}
Up to exchanging the active weights, their pairs are $(2,7)$,
$(3,6)$, and $(4,5)$. For $1\leq r_0\leq7$, none of the corresponding module
dimensions divides $\ell=r_0+8$; the seven checks are listed in
\cref{tab:quartic-nine-low-weights}.
If $r_0\geq8$, the two
smallest weights are the active pair. Hence $a+b-1=8$ and the same
criterion is $8=\delta(8)\mid r_0+8$, equivalently $8\mid r_0$.
\end{proof}

